% Round 10. Agent: g10_trace_audit.
% New deduction, conditional only on the separately audited R6 full finite trace.
% NEEDS_HUMAN_REVIEW=YES. Does not claim general strong G.

\section{Ramified finite trace commutes with restriction as an actual form}
\label{g10t:ramified-base-change}

\paragraph{Status and scope.}
We work over a complete nontrivially valued field $K$ of
characteristic zero. All spaces are analytifications of separated
finite-type $K$-schemes, and all forms and pullbacks below are
Mihatsch's actual $\delta$-forms. The input from the earlier research
is the full finite trace for finite dominant maps between
\emph{normal integral} varieties. This trace accepts arbitrary
actual source forms, is a sheaf map, commutes with $d',d''$, and
satisfies the projection formula and adjunction against every compact
actual test. Its normal, possibly singular, source hypothesis is
essential to the objects below. We do not assume a trace on a
nonreduced fiber product.

The public analytic inputs are Mihatsch, \emph{$\delta$-forms on
Lubin--Tate spaces}, arXiv:2112.10018v4, Definition~4.2,
Theorem~4.5, Lemma~3.14, Corollary~4.9, Proposition~4.11,
and Theorem~5.5. The result proved here strengthens the earlier
top-degree cycle pairing to an equality of actual forms in every
bidegree, including when the testing cycle lies in the branch locus.

\begin{lemma}[Local lifting of actual forms from a closed subspace]
\label{g10t:actual-test-lifting}
Let $i:Z\hookrightarrow X$ be a closed immersion of pure-dimensional
analytic spaces without boundary. Every germ of an actual form on
$Z$ is the pullback of a germ of an actual form on $X$. More
precisely, an expression
\[
                 \eta=t_f^*\gamma,
       \qquad \gamma\in\mathcal B^{a,b}(\mathbf R^s),
\]
near a point $z\in Z$ lifts to the expression $t_F^*\gamma$ near
$i(z)$, with exactly the same Euclidean $\delta$-form $\gamma$.
On paracompact analytic opens, every compactly supported actual
form on $Z$ is a finite sum of restrictions of compactly supported
actual forms on prescribed sufficiently small ambient opens.
\end{lemma}

\begin{proof}
Definition~4.2 supplies the stated presentation. In particular,
$\gamma$ is a $\delta$-form on the ambient Euclidean space; it is
not merely a current balanced on the tropical image of $Z$.
Choose an affinoid neighborhood $U=\mathcal M(A)$ of $i(z)$
small enough that $Z\cap U=\mathcal M(A/I)$ lies in the
presentation open. Each analytic unit $f_j$ on $Z\cap U$ lifts
through $A\twoheadrightarrow A/I$ to an analytic function $F_j$.
Since $F_j(i(z))\ne0$, after shrinking to an ordinary ambient
open all the $F_j$ are units. Thus $t_F^*\gamma$ is a genuine
actual form by Definition~4.2. Theorem~4.5 gives
\[
                    i^*(t_F^*\gamma)=t_f^*\gamma.
\]
This also proves independence of the chosen presentation after
restriction; no new ambient balancing assertion is required.

For the compact statement, cover the compact support by finitely
many such opens. Original-smooth partitions of unity give a finite
decomposition into actual forms whose supports are compactly
contained in the respective restricted opens. Lift each local
piece as above and multiply its lift by an original-smooth ambient
cutoff equal to one near its support on $Z$. The restricted piece
is zero away from that support, so its restriction is unchanged
everywhere on the restricted open. Zero extension of the resulting
compact ambient form gives the assertion. These cutoffs are used
only for extension and are never differentiated in the argument
below.
\end{proof}

\begin{lemma}[Actual forms are determined on a dense algebraic open]
\label{g10t:dense-restriction}
Let $V$ be a pure $m$-dimensional variety and let $S\subset V$
be a closed algebraic subset of dimension at most $m-1$.
The restriction map
\[
        \mathcal B^{p,q}(V^{\rm an})
          \longrightarrow
        \mathcal B^{p,q}((V\setminus S)^{\rm an})
\]
is injective.
\end{lemma}

\begin{proof}
Suppose $\alpha$ restricts to zero. For each compact complementary
original-smooth test $\eta$, the product $\alpha\wedge\eta$ is
an actual $(m,m)$-form. It is zero off $S$, whereas Corollary~4.9
says that its support is disjoint from $S$. Hence the product is
zero. All original-smooth current pairings of $\alpha$ vanish,
and Proposition~4.11 implies $\alpha=0$.
\end{proof}

\begin{lemma}[Top-degree cycle restriction, with scheme lengths]
\label{g10t:cycle-trace}
Let $q:Y\to X$ be finite flat with $X,Y$ smooth of dimension
$n$. Let $i:Z\hookrightarrow X$ be a smooth integral closed
subvariety of dimension $m=n-r$, where $r\geq1$, and put
$W=Y\times_X Z$. For every compact actual
$\beta\in\mathcal B_c^{m,m}(Y^{\rm an})$,
\[
                  \int_Z i^*\operatorname{Tr}_q\beta
                    =\int_W \beta|_W.
\tag{G10.T1}\label{g10t:cycle-trace-eq}
\]
The right side uses the fundamental cycle of the possibly
nonreduced scheme $W$.
\end{lemma}

\begin{proof}
This argument also independently reconstructs the earlier scalar
restriction formula. Work first on an open on which $Z$ is cut
out by a regular sequence $f_1,\ldots,f_r$. Set
\[
 \phi=\min_i v(f_i),\qquad
 g_Z=(-1)^{r-1}\phi(D\phi)^{r-1},\qquad D=d'd''.
\]
Theorem~5.5 gives
\[
 \int g_Z\wedge D\theta=-\int_Z\theta|_Z
\]
for every compact actual $(m,m)$-test $\theta$. Flatness preserves
the regular sequence. The corresponding explicit Green form of
$W$ is therefore $q^*g_Z$, and the same theorem gives
\[
                  \int q^*g_Z\wedge D\beta=-\int_W\beta|_W.
\]
This theorem has no reducedness condition on the complete
intersection. Lemma~3.14 supplies exactly its generic-length
multiplicities in the integration weights.

The support of $D\beta$, of bidegree $(m+1,m+1)$, is disjoint
from $W$ by Corollary~4.9. Hence
$C=q(\operatorname{Supp}D\beta)$ is a compact subset of
$X\setminus Z$. Locality of trace and its commutation with $D$
give
\[
 \operatorname{Supp}(D\operatorname{Tr}_q\beta)
       =\operatorname{Supp}(\operatorname{Tr}_qD\beta)\subseteq C.
\]
Choose an original-smooth compact cutoff $\chi$ on $X\setminus Z$
which is one near $C$. The all-actual trace adjunction applies to
the legitimate compact form $\chi g_Z$, so
\[
 \begin{split}
 \int g_Z\wedge D\operatorname{Tr}_q\beta
  &=\int \chi g_Z\wedge\operatorname{Tr}_q(D\beta)\\
  &=\int q^*(\chi g_Z)\wedge D\beta
   =\int q^*g_Z\wedge D\beta.
 \end{split}
\]
The cutoff is not differentiated. Comparing the two
complete-intersection identities proves
\eqref{g10t:cycle-trace-eq} on this open. A finite
original-smooth partition near $q(\operatorname{Supp}\beta)$,
pulled back to $Y$, reduces the general situation to these opens
and the complement of $Z$. The projection formula then permits
summing the local equalities.
\end{proof}

\begin{theorem}[Form-valued ramified base change along a smooth cycle]
\label{g10t:form-base-change}
Let $q:Y\to X$ be a finite flat surjective map between smooth
integral $n$-dimensional varieties. Let $i:Z\hookrightarrow X$
be a smooth integral closed subvariety of dimension $m$. Write
\[
        [W]=[Y\times_X Z]=\sum_j\ell_j[W_j],
\]
where $W_j$ has its reduced integral structure and
$\ell_j=\operatorname{length}_{\mathcal O_{W,\xi_j}}
\mathcal O_{W,\xi_j}$ is its generic length. Let
$\widetilde W_j\to W_j$ be its normalization, and denote the
induced maps by
\[
             q_j:\widetilde W_j\to Z,
        \qquad \rho_j:\widetilde W_j\to Y.
\]
Then, for every actual $\omega\in\mathcal B^{p,q}(Y^{\rm an})$,
\[
 \boxed{\quad
 i^*\operatorname{Tr}_q\omega
       =\sum_j\ell_j\operatorname{Tr}_{q_j}(\rho_j^*\omega)
       \quad}
\tag{G10.T2}\label{g10t:form-base-change-eq}
\]
as actual forms on $Z^{\rm an}$. The statement is local on
ordinary opens. It permits $Z$ to be contained in the branch
locus and permits every $W_j$ to be singular and $W$ to be
nonreduced.
\end{theorem}

\begin{proof}
If $m=n$, the integral closed subvariety $Z$ is $X$ and the
assertion is tautological; assume $m<n$ in the following argument.
All objects on the right exist before testing. Varieties over
a field have finite normalization. Flatness implies that every
irreducible component of $W$ dominates the integral $Z$, so
each $q_j$ is finite dominant with normal integral source and
normal integral target. Thus the stated finite trace theorem
applies to $q_j$. The map $\rho_j$ is an actual morphism
between pure-dimensional boundaryless analytic spaces, so
Theorem~4.5 defines $\rho_j^*\omega$. There is no invocation
of a trace on the nonreduced scheme $W$.

For a compact actual test
$\eta\in\mathcal B_c^{m-p,m-q}(Z^{\rm an})$, use
Lemma~\ref{g10t:actual-test-lifting} to reduce by a finite sum
to $\eta=i^*\widetilde\eta$ for a compact ambient actual
form $\widetilde\eta$. The form
\[
                  \beta=\omega\wedge q^*\widetilde\eta
\]
is a compact actual $(m,m)$-form on $Y$; compactness follows
from finiteness of $q$. By the projection formula and
Lemma~\ref{g10t:cycle-trace},
\[
 \begin{split}
 \int_Z (i^*\operatorname{Tr}_q\omega)\wedge\eta
   &=\int_Z i^*\operatorname{Tr}_q
                        (\omega\wedge q^*\widetilde\eta)\\
   &=\int_W(\omega\wedge q^*\widetilde\eta)|_W\\
   &=\sum_j\ell_j\int_{\widetilde W_j}
                        \rho_j^*\omega\wedge q_j^*\eta\\
   &=\int_Z\left(\sum_j\ell_j
                  \operatorname{Tr}_{q_j}(\rho_j^*\omega)\right)
                         \wedge\eta.
 \end{split}
\]
For the normalization step, Lemma~3.14 first decomposes the
integral on $W$ with the coefficients $\ell_j$. An actual
top-degree form on $W_j$ is supported away from any proper
algebraic closed subset by Corollary~4.9. We may therefore
remove the locus where normalization is not an isomorphism,
without changing either integral. This proves that step
without a trace theorem for a nonnormal target.

All compact actual tests have now been covered. In particular
all original-smooth tests have been covered, and actual-current
injectivity on $Z$ proves the claimed equality of forms.
If a complementary degree is negative, both sides vanish by
the dimension bound for actual forms, so there is no omitted
bidegree.
\end{proof}

\begin{corollary}[Normal testing varieties and normalization of cycles]
\label{g10t:normal-cycle-base-change}
The same equality holds when $Z$ is normal integral, with
$i:Z\hookrightarrow X$ a closed immersion. More generally,
if $Z_0\subset X$ is any integral closed subvariety and
$h:\overline Z_0\to X$ is the composite of its normalization
with its inclusion, then the equality holds with $i$ replaced
by $h$, using the generic-length components of
$Y\times_X\overline Z_0$.
\end{corollary}

\begin{proof}
In the first case compare both actual forms on the smooth
locus of $Z$. In the second, compare them on the dense open
of $\overline Z_0$ identified with the smooth locus of $Z_0$.
The previous theorem applies there, since that smooth locus
is a closed subvariety of $X\setminus (Z_0)_{\rm sing}$.
The finite traces on the normal components exist globally,
and they commute with restriction to this open. Both sides
are therefore actual forms on the normal testing variety
which agree on a dense algebraic open. Apply
Lemma~\ref{g10t:dense-restriction}.
\end{proof}

\paragraph{Ramification check.}
For $q:\mathbf A^1_t\times Z\to\mathbf A^1_x\times Z$,
$x=t^e$, and the zero section $i:Z\hookrightarrow
\mathbf A^1_x\times Z$, the fiber product is
$\operatorname{Spec}(\mathcal O_Z[t]/(t^e))$. It has one
reduced component, normalization $Z$, and generic length $e$.
Formula~\eqref{g10t:form-base-change-eq} becomes
\[
                i^*\operatorname{Tr}_q\omega
                        =e\,\omega|_{\{t=0\}\times Z}.
\]
Keeping only reduced components with coefficient one would
already fail on the constant actual form $\omega=1$.

\paragraph{Consequence for the strong-G program.}
This is a compatibility result for arbitrary actual coefficients,
all test forms, original-field degrees, and cycle restriction.
It supplies the form-valued branch-inclusive interface that a
top-degree scalar projection formula alone does not supply.
It does not construct a zero-curvature Green form for an
arbitrary homologically trivial cycle, and it does not give
an actual pushforward for a map with positive-dimensional
fibers. Its only nonpublic theorem input remains the separately
audited full finite actual trace.
